\documentclass[11pt,a4paper]{article}

% --- packages (all standard in texlive) ---
\usepackage[utf8]{inputenc}
\usepackage[T1]{fontenc}
\usepackage{lmodern}
\usepackage[margin=2.4cm]{geometry}
\usepackage{amsmath,amssymb}
\usepackage{booktabs}
\usepackage{array}
\usepackage{longtable}
\usepackage{enumitem}
\usepackage{listings}
\usepackage{xcolor}
\usepackage{parskip}
\usepackage[hidelinks]{hyperref}

% --- light styling (matches the connectivity-estimation chapter) ---
\definecolor{accent}{RGB}{40,70,140}
\definecolor{codebg}{RGB}{246,247,250}
\definecolor{codekw}{RGB}{40,70,140}
\definecolor{codecm}{RGB}{110,120,130}
\hypersetup{colorlinks=true,linkcolor=accent,urlcolor=accent,citecolor=accent}

\lstset{
  basicstyle=\ttfamily\small,
  backgroundcolor=\color{codebg},
  keywordstyle=\color{codekw}\bfseries,
  commentstyle=\color{codecm}\itshape,
  stringstyle=\color{accent},
  showstringspaces=false,
  breaklines=true,
  frame=single,
  rulecolor=\color{codebg},
  framesep=6pt,
  xleftmargin=4pt,
  columns=fullflexible,
  keepspaces=true,
}

\newcommand{\code}[1]{\texttt{#1}}
\newcommand{\config}{\code{ExperimentConfig}}

\title{\textbf{A Reproducible, Scalable Framework for\\ Connectivity-Inference Experiments}\\[4pt]
\large From a Laptop to a Cluster, and from Private Code to a FAIR Testbed}
\author{Mohammad Joudy \\ \small Bernstein Center Freiburg}
\date{\today}

\begin{document}
\maketitle

\begin{abstract}
\noindent
The previous chapters developed a \emph{method}: from calcium fluorescence, through
deconvolution and preprocessing, to an estimate of synaptic connectivity. A method,
however, is only as valuable as our ability to \emph{reproduce} it, \emph{scale} it, and
let \emph{others} build on it. This closing chapter is about the instrument that carries
the method: a software framework that runs the entire pipeline --- network simulation,
calcium generation, preprocessing, and connectivity inference --- as a single,
configuration-driven experiment that behaves identically on a laptop and on a
high-performance-computing (HPC) cluster, records exactly which code and which parameters
produced every number, and is designed from the ground up to be extended by people other
than its author. We first state the problem (why personal research code rarely survives its
researcher), then describe the architecture and the concrete reproducibility mechanisms,
then show how the same code moves from a small $N=100$ network on a laptop to the full
Brunel network ($N=12500$) on a GPU, and finally recast the whole system through the FAIR
principles \cite{wilkinson2016} as a ground-truth \emph{testbench} for connectivity
inference. No experimental results are re-presented here --- those belong to the earlier
chapters; the object of study in this chapter is the experimental apparatus itself.
\end{abstract}

\tableofcontents
\vspace{1em}
\hrule
\vspace{1em}

% ======================================================================
\section{Why a chapter about the software}
\label{sec:why}

A computational-neuroscience thesis usually reports a result: a network was simulated, a
signal was analysed, a conclusion was drawn. The code that produced the result is treated as
scaffolding --- necessary to build the argument, then quietly discarded. This chapter takes
the opposite view. It argues that, for a study like this one, \textbf{the reproducible
experimental apparatus is itself a scientific contribution}, and it documents that apparatus
as carefully as the earlier chapters documented the method.

There are two reasons to do this, one defensive and one constructive.

\paragraph{The defensive reason: results that cannot be reproduced are not yet results.}
Connectivity inference from calcium imaging is a chain of many steps --- a stochastic network
simulation, a nonlinear calcium model, a preprocessing pipeline with several tunable windows,
a regularised regression with a delicate hyper-parameter, and a post-processing stage. Each
step has parameters; several have hidden state (random seeds, thread counts, library
versions). A single wrong lag or a silently different random stream can move a correlation
from meaningful to zero. If, six months later, we cannot say \emph{exactly} which code and
which configuration produced a given figure, we cannot defend it, extend it, or trust it. The
framework exists so that every number in this thesis is traceable to an exact line of code, an
exact configuration, and an exact software environment.

\paragraph{The constructive reason: the method should outlive the researcher.}
The pipeline of the earlier chapters was, at first, a collection of personal scripts written to
answer one person's questions on one network. That is the normal state of research code, and
it is also the reason so little of it is ever reused. The worry that motivates this chapter is
concrete: \emph{if someone --- a labmate, a reviewer, a future student --- wanted to test one
of my ideas, change one assumption, or add a competing method, could they?} Turning ``my code''
into ``a resource'' is not a matter of tidiness; it is a matter of \emph{generality} ---
whether the software lets a stranger ask a \emph{different} question than the one it was born to
answer. The design decisions in this chapter are the ones that make that possible.

\paragraph{What makes this a neuroscience contribution, not only an engineering one.}
The decisive fact is that the networks here are \emph{simulated}, so the true connectivity is
\textbf{known}. Real calcium-imaging experiments never provide this: one infers connectivity and
can almost never check the answer against the truth. A framework that generates ground-truth
networks, degrades them into realistic calcium signals under controlled conditions, and runs
identically from a laptop to a cluster is therefore not merely reproducible plumbing --- it is a
\textbf{testbench} in which any connectivity-inference method can be measured against a known
answer, across network sizes and dynamical regimes, by anyone. That is the contribution this
chapter claims.

% ======================================================================
\section{What ``reproducible'' has to mean here}
\label{sec:repro-defn}

``Reproducible'' is often used loosely. For this framework it has a precise, testable meaning,
which we can state as an equation of intent:
\[
\text{same code} \;+\; \text{same configuration} \;+\; \text{same seeds}
\;\;\Longrightarrow\;\; \text{same result, on any machine.}
\]
Every word on the left is something the framework must \emph{pin down} and \emph{record}:

\begin{description}[leftmargin=1.4em,style=nextline]
  \item[Same code.] Not ``the same repository'', but the same \emph{commit} --- a specific,
    immutable snapshot of every source file. The framework stamps the git commit hash (and a
    flag marking any uncommitted changes) into every run, so a result can never drift away from
    the code that made it.
  \item[Same configuration.] Every parameter of every stage --- network size, connection
    density, calcium time constant, smoothing window, regression lag, regularisation strength,
    solver, random seed --- lives in one place, a single configuration object, and is saved
    verbatim alongside the run.
  \item[Same seeds.] The simulation is stochastic. Reproducibility requires fixing the random
    seed \emph{and} --- an honest subtlety documented in the framework --- the number of
    simulation threads, because the parallel random-number stream depends on both. Fixing one
    without the other reproduces the \emph{style} of a result but not the exact numbers.
  \item[On any machine.] The software environment (every library, at an exact version) must be
    reconstructible on a laptop and on a cluster alike. This is the job of a pinned environment
    and, for the durable record, a container.
\end{description}

The rest of the chapter is, in effect, a walk through how each of these four requirements is
met in code.

% ======================================================================
\section{Architecture: a configuration-driven pipeline}
\label{sec:arch}

The software is an installable Python package, \code{calcium\_ar}, organised so that each
stage of the scientific pipeline is a separate, testable module and the stages are wired
together by a single experiment runner.

\paragraph{The four stages.}
The pipeline mirrors the physical story of the thesis:
\begin{center}
\code{SIMULATION} $\;\longrightarrow\;$ \code{CALCIUM} $\;\longrightarrow\;$
\code{PREPROCESSING} $\;\longrightarrow\;$ \code{SOLVER (+POST-PROCESS)}
\end{center}
\begin{enumerate}[leftmargin=1.6em]
  \item \textbf{Simulation} (\code{simulation/}) builds and runs a Brunel-type balanced
    random network in NEST \cite{nest}, returning binary spike trains and the true adjacency
    matrix --- the ground truth.
  \item \textbf{Calcium} (\code{simulation/calcium\_signal.py}) turns spikes into fluorescence
    traces through the AR-style calcium model of the earlier chapters.
  \item \textbf{Preprocessing} (\code{preprocessing/}) recovers a per-neuron activity feed
    (smoothing, derivative, spike-cut, robust time-constant estimation, feed reconstruction).
  \item \textbf{Solvers} (\code{solvers/}) fit the connectivity by regression, with a family
    of interchangeable backends and an optional post-processing stage.
\end{enumerate}

\paragraph{The central idea: one configuration object, one runner.}
Everything a run needs is described by a single dataclass, \config, and executed by a single
function, \code{run\_single(config)}, which returns a structured \code{ExperimentResult}
(metrics plus the paths to the artefacts it produced). This one decision is what makes the
whole system reproducible \emph{and} scalable: there is exactly one place where the meaning of a
run is defined, and it is plain data that can be saved, compared, and regenerated.

\begin{lstlisting}[language=Python,caption={An experiment is a value, not a script. The same call runs on a laptop and on a cluster node.},captionpos=b]
from calcinet.framework import ExperimentConfig, run_single

cfg = ExperimentConfig(
    n_excitatory=80, n_inhibitory=20,   # the network (regime)
    epsilon=0.1, g=5.0, eta=2.0,        # dynamical regime
    tau=100.0, sigma_extra=0.0,         # the calcium signal
    lag_ms=1.5, solver="fista",         # the estimator
    lam=3e-3, lam_l2=1e-3,
    seed=1, n_threads=8,                # reproducibility knobs
    device="cpu",                       # "cuda" on the GPU path
)
result = run_single(cfg)                # -> ExperimentResult (+ a run directory)
\end{lstlisting}

\paragraph{Sweeps are the same idea, lifted.}
A scientific question is rarely a single run; it is usually ``how does the result change as I
vary this knob?''. The framework expresses that directly: a \code{ParameterSweep} takes a base
configuration and a dictionary of parameters to vary, and expands into the full set of
\config\ objects, each of which is just another \code{run\_single}. Because a sweep is nothing
more than a list of configurations, the \emph{same} abstraction that runs one experiment on a
laptop becomes, unchanged, a job array on a cluster (Section~\ref{sec:hpc}). Results are
collected back into a single table by \code{load\_sweep\_results()}, which returns one row per
run with every configuration field and every metric.

\paragraph{Two kinds of knobs.}
It is worth naming a distinction that the configuration object makes visible and that organises
the entire experimental programme. The parameters split into two roles with fundamentally
different meanings:
\begin{itemize}[leftmargin=1.6em]
  \item \textbf{Estimator knobs} --- how we \emph{recover} connectivity from a fixed dataset
    (lag, solver, regularisation, post-processing). These are what one \emph{tunes}. They were
    the subject of the previous chapter and are, on a small network, essentially settled.
  \item \textbf{Regime knobs} --- what network and signal we are \emph{trying} to recover
    (size $N$, density, the excitation/inhibition ratio $g$, firing rate, calcium time constant,
    recording length, noise). These \emph{define the problem}. Held fixed while the method was
    developed, they become the \emph{axes of generalisation} that the HPC phase exists to
    explore.
\end{itemize}
This is not merely bookkeeping: it tells us which parameters are answers and which are
questions, and it is the reason the framework is organised so that changing a regime knob is a
one-line edit rather than a rewrite.

\paragraph{A solver family built for scale.}
Connectivity inference is a regression whose design matrix grows with both the number of
neurons $N$ and the recording length $T$. Written naively, the normal equations hold an
$O(N\cdot T)$ matrix in memory and do not fit for large networks. The \code{solvers/} module
therefore provides a family of interchangeable backends behind a common \code{solve\_*()}
interface --- exact normal equations, ridge, and several PyTorch backends (a GPU normal-equation
solver, a mini-batch $L_1$/elastic-net solver, and a linear-layer solver) --- all of which
process the data in \emph{chunks} so that memory stays $O(N^2)$ rather than $O(N\cdot T)$. The
common interface matters twice over: it is what lets a large run move from CPU to GPU by
changing one field (\code{device="cuda"}), and it is the extension point at which a newcomer
plugs in \emph{their} method and inherits the entire reproducible harness for free
(Section~\ref{sec:testbed}).

% ======================================================================
\section{The reproducibility mechanisms}
\label{sec:mechanisms}

Section~\ref{sec:repro-defn} listed four things that must be pinned down. Here is how each is
realised.

\subsection{Deterministic simulation (same seeds)}
The network simulation is stochastic, so determinism is arranged deliberately. The random seed
is a first-class field of \config\ and is passed into the NEST kernel; with the seed fixed, the
same configuration yields the same network and the same spike trains. The framework documents an
important and easily-missed caveat: the parallel random-number stream also depends on the number
of simulation threads, so \textbf{exact} reproducibility requires fixing the seed \emph{and}
\code{n\_threads}. This is exactly the kind of hidden state that silently breaks reproductions,
and the framework's response is to make it explicit rather than to hide it --- an honest record
of a real constraint is worth more than a false promise of push-button determinism.

\subsection{Pinned environments and a container (any machine)}
Two levels of environment control are provided, matching two levels of durability.
\begin{itemize}[leftmargin=1.6em]
  \item A curated \code{environment.yml} together with an exact \code{environment.lock.yml}
    reconstructs the software stack down to the version of every transitive dependency. This is
    the low-friction path and is what runs day to day.
  \item For the archival record, the framework's deployment design specifies an
    \textbf{Apptainer} (Singularity) container rather than Docker, because HPC systems do not
    grant the root privileges Docker assumes. The image is built once and becomes the single
    immutable artefact that reproduces the environment years later; the container definition and
    the lockfile live in version control, while the built image is archived alongside the data.
\end{itemize}
The principle is that the environment is not a description in prose (``we used Python~3.11 and
recent versions of the usual libraries'') but a reconstructible object.

\subsection{Provenance stamping (same code)}
This is the mechanism that ties results to code. On every run, the framework queries the git
revision --- the commit hash plus a ``dirty'' flag that records whether the working tree had
uncommitted edits --- and writes it into both the run's own directory and the central ledger.
The consequence is a hard guarantee: \emph{no number in this thesis exists without a stored
pointer to the exact code that produced it}. A dirty flag on a run is a visible warning that its
code was not fully committed, and therefore not fully reproducible --- the framework would rather
flag that than let it pass unnoticed.

\subsection{The experiment record system (two-layer memory)}
Reproducibility is not only technical; it is also a matter of remembering what was done and why.
The framework keeps two complementary records:
\begin{itemize}[leftmargin=1.6em]
  \item A \textbf{ledger} (\code{results/ledger.csv}): one row per run, carrying every
    configuration field, every metric, the run directory, the timestamp, and the git commit.
    It is appended automatically by the runner and queried by a small command-line tool
    (\code{src/calcinet/ledger_cli.py}), which can, for example, tabulate the effect of a single
    parameter on a chosen metric. The ledger is the machine-readable, searchable index of the
    entire investigation.
  \item A \textbf{notebook} (\code{docs/experiments/notebook.md}): dated, human-readable entries
    that record the \emph{conclusions} drawn from each campaign. Where the ledger holds the
    numbers, the notebook holds the meaning.
\end{itemize}
Together they are a durable memory of the project: the numbers, and the reasons.

\subsection{Git as the only transport}
Finally, the framework fixes \emph{how code and results move} between machines. There is one
repository and three locations --- the laptop where code is written, a hosted origin, and the
cluster \code{\$HOME} that only ever pulls and runs. Code flows in one direction only
(laptop~$\rightarrow$~origin~$\rightarrow$~cluster); deployments are made from an immutable git
\emph{tag}, never a moving branch, so a cluster run is anchored to a snapshot that cannot change
underneath it. Results flow back the other way: text (the ledger, configurations) through git
because it is small and diff-able, and large arrays through \code{rsync} from the cluster
workspace, because they do not belong in version control. Code is never hand-edited on the
cluster without being committed and pushed first. This one-way discipline is unglamorous, and it
is the difference between a reproducible pipeline and a folder of files nobody can account for.

% ======================================================================
\section{One code, two scales: from a laptop to the cluster}
\label{sec:hpc}

The payoff of the design is that moving from a small exploratory network to a large one is a
change of \emph{configuration and compute backend}, not a change of code.

\paragraph{A dual compute path, single node.}
Simulation and preprocessing are cheap and run on the CPU (NumPy, SciPy, scikit-learn,
NEST). The regression at large $N$ is the expensive part and runs on a GPU through the PyTorch
solver backends. The framework deliberately stays \textbf{single-node}: NEST parallelises with
threads rather than across nodes, which avoids the considerable pain of binding a multi-node MPI
job inside a container, and the GPU path needs only the container's \code{--nv} flag to see the
device. Simplicity here is a feature --- it is what keeps the containerised, reproducible run
identical to the laptop run.

\paragraph{Sweeps become job arrays.}
Because a parameter sweep is just a list of configurations
(Section~\ref{sec:arch}), it maps directly onto a SLURM \emph{job array}: each array index runs
one \code{run\_single} with one configuration, writes its own run directory, and appends its own
ledger row. The same \code{ParameterSweep} object that a reader can run on a laptop over lunch is,
without modification, the specification of a thousand-run campaign on the cluster. Nothing about
the experiment's \emph{meaning} changes between the two; only the executor does.

\paragraph{The scaling target: the full Brunel network on a GPU.}
The concrete horizon this framework is built for is the \textbf{full Brunel network of
$N=12500$ neurons} run on a GPU. Reaching it is not a rewrite but a coordinated turn of a few
knobs the framework already exposes:
\begin{itemize}[leftmargin=1.6em]
  \item the regime knobs $N$, density, and the input strength are set to the target network,
    respecting the scaling rule that keeps the network in the correct balanced regime as $N$
    grows (the per-connection weight must be rescaled with the number of inputs, or the dynamics
    move to the wrong regime and the inference becomes meaningless);
  \item the recording length is budgeted upward, because a larger network has more parameters to
    estimate per neuron and therefore needs more data;
  \item the compute path switches to \code{device="cuda"} with a chunk size chosen to fit GPU
    memory, which the chunked solvers were written precisely to allow.
\end{itemize}
The scientific reason this scale matters --- whether the method's accuracy survives the greater
collinearity of a large, densely-recorded network --- is a question for the experiments, not for
this chapter. The point here is architectural: the framework is designed so that this question
can be \emph{asked} by editing a configuration, and so that its answer, whatever it is, arrives
with full provenance and is reproducible by anyone with the same container.

% ======================================================================
\section{FAIR by design}
\label{sec:fair}

The reproducibility mechanisms above were introduced pragmatically, one problem at a time. Seen
together, they are an implementation of a recognised set of principles for research objects: the
\textbf{FAIR} principles --- Findable, Accessible, Interoperable, Reusable
\cite{wilkinson2016}. FAIR was formulated for scientific \emph{data}, but its authors are
explicit that it applies equally to the non-data research objects --- the algorithms, tools, and
workflows --- that produce and consume that data. A connectivity-inference pipeline is exactly
such a workflow, and it is illuminating to read the framework back through the four principles.

Table~\ref{tab:fair} maps each FAIR sub-principle to the concrete mechanism that serves it.

\begin{longtable}{@{}p{0.30\linewidth} p{0.64\linewidth}@{}}
\caption{The framework's mechanisms mapped onto the FAIR principles \cite{wilkinson2016}.}
\label{tab:fair}\\
\toprule
\textbf{FAIR sub-principle} & \textbf{How the framework realises it}\\
\midrule
\endfirsthead
\toprule
\textbf{FAIR sub-principle} & \textbf{How the framework realises it}\\
\midrule
\endhead
\multicolumn{2}{@{}l}{\textit{Findable}}\\
F1 --- persistent identifier & Each run is a directory keyed by its configuration and timestamp;
each deployable version is an immutable git tag; the archival endpoint is a released,
DOI-minted snapshot (the last step of the roadmap, Section~\ref{sec:status}).\\
F2, F3 --- rich metadata, bound to the data & Every run stores its full \config\ and its git
commit next to the arrays it produced; metadata and data are inseparable by construction.\\
F4 --- indexed, searchable & The ledger (\code{ledger.csv}) is the searchable index of every
run; \code{src/calcinet/ledger_cli.py} queries it by parameter and metric.\\
\midrule
\multicolumn{2}{@{}l}{\textit{Accessible}}\\
A1 --- retrievable by open protocol & Code and text move over git; large arrays over rsync ---
open, universally available protocols.\\
A2 --- metadata persist beyond the data & The ledger and configurations are text under version
control; even when a large array is pruned to save space, its provenance row survives.\\
\midrule
\multicolumn{2}{@{}l}{\textit{Interoperable}}\\
I1 --- open, shared formats & Plain, non-proprietary formats throughout: CSV, JSON, NumPy/Zarr
arrays, a documented Python API.\\
I2 --- a shared vocabulary & The \config\ schema is a single controlled vocabulary of parameter
names used identically by the runner, the ledger, the sweeps, and the notebook.\\
I3 --- qualified references & Every result references the exact commit and dataset it derives
from, so provenance links are typed, not implied.\\
\midrule
\multicolumn{2}{@{}l}{\textit{Reusable}}\\
R1 --- rich attributes & Each run carries its complete configuration and its complete metrics.\\
R1.1 --- a clear licence & An OSI-approved licence is the small remaining step toward full reuse
(Section~\ref{sec:status}).\\
R1.2 --- detailed provenance & Git commit and dirty flag, the environment lockfile, and the
container together constitute complete, machine-checkable provenance.\\
R1.3 --- community standards & The framework is built on the field's standard tools --- the
Brunel model and NEST, scikit-learn and PyTorch, conda, SLURM, and Apptainer --- so its outputs
are legible to the community that would reuse them.\\
\bottomrule
\end{longtable}

The value of reading the framework this way is not to collect compliance badges. It is that FAIR
gives a precise language for the constructive worry of Section~\ref{sec:why}: the distance
between ``code that works for me'' and ``a resource others can use'' is exactly the distance
across this table, and the table shows that the distance has largely been closed by decisions
made for reproducibility's own sake.

% ======================================================================
\section{A ground-truth testbench for the community}
\label{sec:testbed}

Reproducibility and FAIRness are means; the end is a \textbf{testbench}. Because the framework
simulates its networks, it holds the one thing that real experiments cannot supply: the true
connectivity matrix. This turns the software from a private pipeline into a controlled arena in
which claims about connectivity inference can be \emph{measured against a known answer}.

\paragraph{The extension points.}
Three seams make the arena open to others:
\begin{itemize}[leftmargin=1.6em]
  \item A new \textbf{inference method} is a new \code{solve\_*()} behind the common solver
    interface; it immediately inherits the whole harness --- provenance, ledger, sweeps, HPC
    execution --- and can be scored against every existing method on identical data.
  \item A new \textbf{signal or preprocessing model} slots into the corresponding stage without
    disturbing the others, because the stages communicate through plain arrays.
  \item A new \textbf{regime} --- a different size, density, excitation/inhibition ratio, firing
    rate, or calcium time constant --- is a change of regime knobs, i.e.\ a different point in the
    space the framework was built to sweep.
\end{itemize}

\paragraph{What this enables.}
Concretely, the framework opens three roads. First, the \emph{generalisation study}: sweeping
the regime knobs one axis at a time to ask whether the method of the earlier chapters, settled on
a small network, holds as the problem changes. Second, \emph{scaling} to the full Brunel network
on GPU (Section~\ref{sec:hpc}), where the interesting scientific risk --- does the estimator's
accuracy survive the collinearity of a large network? --- can finally be tested. Third, and most
importantly for the argument of this chapter, \emph{reuse by others}: a labmate, a reviewer, or a
future student can take the container, reproduce any result exactly, drop in a competing method,
and get a comparable number --- which is the operational definition of the method having outlived
its author.

\subsection{Current status and the road to full FAIRness}
\label{sec:status}
In the spirit of the FAIR principles --- which are meant to be adopted incrementally, not all at
once --- it is worth stating plainly what is in place and what remains. Firmly in place: the
installable package and configuration-driven runner; deterministic simulation with documented
seed/thread handling; pinned and lockfile-exact environments; git provenance stamping on every
run; the ledger-plus-notebook record system; the git-as-transport workflow; the chunked CPU/GPU
solver family; and the SLURM job-array scaffolding used for the first larger-network campaigns.
The remaining steps, each small and each explicitly a FAIR item from Table~\ref{tab:fair}, are:
adding an OSI-approved licence (R1.1), publishing a versioned, DOI-minted release of the code and
a reference dataset (F1), and building and archiving the Apptainer image as the citable
environment artefact (R1.2). None of these changes the science; each moves the framework further
along the continuum from a personal instrument to a community resource.

% ======================================================================
\section{Summary}
\label{sec:summary}

This chapter documented the instrument behind the thesis rather than a new result. Its argument
was that, for a multi-stage computational study whose credibility rests on many interacting
parameters and hidden stochastic state, a reproducible experimental framework is a genuine
contribution --- and, because the networks are simulated and their connectivity known, a
scientifically useful one. The framework realises a precise notion of reproducibility (same code,
configuration, and seeds give the same result on any machine) through concrete mechanisms:
a single configuration object and runner, provenance stamping against exact git commits, pinned
environments and a container, a two-layer ledger-and-notebook record, and a one-way
git-as-transport discipline. The same abstractions that make one run reproducible on a laptop
make a thousand-run sweep reproducible on a cluster, and carry the pipeline, by a change of
configuration rather than of code, toward the full Brunel network of $N=12500$ neurons on a GPU.
Read through the FAIR principles, these mechanisms are most of what separates private research
code from a reusable, ground-truth testbench that others can extend --- which is the lasting
contribution this final chapter claims for the software behind the method.

% ======================================================================
\begin{thebibliography}{9}
\bibitem{wilkinson2016}
M.~D. Wilkinson \emph{et al.},
``The FAIR Guiding Principles for scientific data management and stewardship,''
\emph{Scientific Data}, \textbf{3}:160018, 2016. doi:10.1038/sdata.2016.18.

\bibitem{nest}
M.-O. Gewaltig and M.~Diesmann,
``NEST (NEural Simulation Tool),'' \emph{Scholarpedia}, \textbf{2}(4):1430, 2007.

\bibitem{brunel2000}
N.~Brunel,
``Dynamics of sparsely connected networks of excitatory and inhibitory spiking neurons,''
\emph{Journal of Computational Neuroscience}, \textbf{8}(3):183--208, 2000.

\bibitem{sklearn}
F.~Pedregosa \emph{et al.},
``Scikit-learn: Machine Learning in Python,''
\emph{Journal of Machine Learning Research}, \textbf{12}:2825--2830, 2011.

\bibitem{pytorch}
A.~Paszke \emph{et al.},
``PyTorch: An Imperative Style, High-Performance Deep Learning Library,''
\emph{Advances in Neural Information Processing Systems}, \textbf{32}, 2019.

\bibitem{kurtzer2017}
G.~M. Kurtzer, V.~Sochat, and M.~W. Bauer,
``Singularity: Scientific containers for mobility of compute,''
\emph{PLOS ONE}, \textbf{12}(5):e0177459, 2017.
\end{thebibliography}

\end{document}
